\label{sec:Marco-referencia}


%--------------------------------------------------------------------------------------
%--------------------------------------------------------------------------------------
%--------------------------------------------------------------------------------------
\section{Marco teórico}
\label{sec:Marco-teorico}
    El estudio de la pobreza desde una perspectiva territorial exige bases conceptuales y metodológicas que permitan comprender tanto la naturaleza multidimensional del fenómeno como los desafíos asociados a su medición en escalas geográficas pequeñas. En particular, la combinación de encuestas de hogares con fuentes censales y registros administrativos abre la posibilidad de utilizar metodologías estadísticas avanzadas que compensen la falta de representatividad muestral en dominios territoriales pequeños.  
    
    Este capítulo presenta los fundamentos teóricos que sustentan el uso de técnicas de SAE, sus extensiones espaciales y los aportes recientes de la ciencia de datos para el análisis territorial de la pobreza. Se integran además antecedentes relevantes que demuestran el uso y la pertinencia de estas metodologías en América Latina y Colombia. Las definiciones, ecuaciones y criterios que se presentan a continuación son los que se aplican operativamente en la Sección~\ref{sec:metodologia} y cuyos valores se reportan en la Sección~\ref{sec:resultados}.

%--------------------------------------------------------------------------------------
%--------------------------------------------------------------------------------------
\subsection{Contexto del estudio}

%--------------------------------------------------------------------------------------
\subsubsection{Índice de Pobreza Multidimensional (IPM)}

La pobreza multidimensional reconoce que las privaciones que enfrentan los hogares trascienden el ingreso e incluyen dimensiones como salud, educación, vivienda, servicios básicos y empleo. Su medición se apoya en el método de conteo propuesto por Alkire y Foster \cite{alkire2011}, que identifica a los hogares privados en un número suficiente de indicadores y agrega esa información en una medida única. En Colombia, la medición oficial se realiza mediante el IPM, lo que permite identificar privaciones específicas y orientar intervenciones focalizadas \cite{DANE_PobrezaMultidimensional_2023}.

El método de Alkire y Foster parte de una matriz de privaciones en la que cada hogar $h$ recibe un valor $g_{hj}$ igual a uno si presenta privación en el indicador $j$ y cero en caso contrario. A cada indicador se le asigna un peso $w_j$, de modo que el puntaje de privación del hogar corresponde a la suma ponderada de sus privaciones, según la Ecuación~\ref{eq:puntaje-privacion}.

\begin{equation}
    c_h = \sum_{j=1}^{J} w_j\, g_{hj},
    \qquad \text{con} \quad \sum_{j=1}^{J} w_j = 1
    \label{eq:puntaje-privacion}
\end{equation}

\noindent donde:

\begin{itemize}
    \item $c_h$: puntaje de privación del hogar $h$.
    \item $g_{hj}$: indicador binario de privación del hogar $h$ en el indicador $j$.
    \item $w_j$: peso asignado al indicador $j$.
    \item $J$: número total de indicadores considerados.
\end{itemize}

Un hogar se identifica como pobre multidimensional cuando su puntaje de privación alcanza o supera un umbral $k$, fijado en la medición oficial colombiana en un tercio de las privaciones ponderadas. La incidencia de pobreza multidimensional del dominio corresponde entonces a la proporción de personas que residen en hogares identificados como pobres, tal como expresa la Ecuación~\ref{eq:incidencia-ipm}.

\begin{equation}
    H = \frac{1}{N} \sum_{h=1}^{N} \mathbb{1}(c_h \geq k)
    \label{eq:incidencia-ipm}
\end{equation}

\noindent donde $N$ es el total de personas del dominio, $\mathbb{1}(\cdot)$ es la función indicadora y $k$ el umbral de identificación. Esta incidencia es la cantidad que el presente trabajo estima a escala municipal y la que, calculada sobre la ECV, constituye la estimación directa departamental descrita en la Subsección~\ref{subsec:estimadores-directos}.

La disponibilidad territorial del IPM es limitada: salvo la estimación derivada del Censo 2018, no existen mediciones periódicas para los 1{,}122 municipios del país \cite{DANE_PobrezaMultidimensional_2023}. Esta brecha de información dificulta el diagnóstico local y la formulación de políticas basadas en evidencia, especialmente en zonas con baja representatividad estadística, y motiva el uso de técnicas capaces de producir estimaciones confiables a nivel municipal mediante fuentes auxiliares, como SAE \cite{DANE_ProyectoSAE_2023}.

%--------------------------------------------------------------------------------------
\subsubsection{Modelos SAE}

Los modelos SAE ofrecen una base estadística sólida para generar estimaciones precisas en dominios con escasa muestra, mediante la integración de información proveniente de censos, registros administrativos y encuestas \cite{CEPAL_PovertyMapping_2022}. Existen dos grandes enfoques: los modelos de área, como el modelo Fay-Herriot, que utilizan estimaciones directas agregadas y covariables auxiliares; y los modelos de unidad, que emplean microdatos para capturar variabilidad individual \cite{RaoMolina_2015}. La selección depende de la calidad y granularidad de las fuentes disponibles, pero ambos enfoques permiten mejorar la precisión, reducir la varianza y generar estimaciones territorialmente consistentes, lo cual es fundamental en contextos municipales.

No obstante, cuando los indicadores presentan patrones territoriales no explicados completamente por las covariables, resulta necesario recurrir a enfoques espaciales que modelen explícitamente la dependencia geográfica. Los modelos espaciales integran estructuras de vecindad o efectos aleatorios espaciales, lo que permite suavizar las estimaciones entre municipios contiguos y reducir variabilidad no explicada. La literatura demuestra que estos modelos mejoran la coherencia territorial, especialmente en dominios pequeños y heterogéneos como los municipios colombianos \cite{PratesiSalvati_2019}.

Las herramientas de ciencia de datos complementan la metodología de estimación en tres frentes. El primero es la selección de variables auxiliares, apoyada en criterios de relevancia temática, disponibilidad de fuentes oficiales y diagnósticos econométricos como el Factor de Inflación de Varianza descrito en la Subsección~\ref{subsec:ols}, que permite descartar predictores redundantes \cite{Greene_2018}. El segundo es la incorporación de fuentes heterogéneas, entre ellas información satelital como el Visible Infrared Imaging Radiometer Suite (VIIRS), cuya luminosidad nocturna opera como aproximación de la actividad económica municipal. El tercero es la visualización mediante mapas temáticos, que facilita a instituciones y tomadores de decisión la interpretación de los patrones territoriales resultantes \cite{CEPAL_DPS_MapaPobreza_2021}.

La aplicación de modelos SAE y de sus extensiones espaciales depende directamente de la disponibilidad de covariables confiables y de cobertura completa. En este proyecto se emplean fuentes oficiales como el Censo Nacional de Población y Vivienda 2018, la GEIH y la ECV del DANE, registros administrativos sectoriales y datos geoespaciales abiertos \cite{DANE_PobrezaMultidimensional_2023, DANE_ENCV_2022, dane2024}. A partir de estas fuentes se derivan variables relacionadas con condiciones de vivienda, servicios públicos, educación, entorno e infraestructura, orientadas a mejorar la capacidad de estimación del modelo y compensar la baja representatividad muestral en muchos municipios.

Estas herramientas conceptuales y metodológicas sustentan el desarrollo del proyecto al permitir generar estimaciones municipales del IPM más precisas y territorialmente coherentes, y aportan una base para el análisis territorial en contextos con limitaciones de información estadística.

%--------------------------------------------------------------------------------------
%--------------------------------------------------------------------------------------
\subsection{Métodos y modelos}

%--------------------------------------------------------------------------------------
\subsubsection{Estimadores directos}
\label{subsec:estimadores-directos}

Se denomina estimador directo a aquel que utiliza exclusivamente la información muestral del dominio de interés, sin recurrir a información externa ni a supuestos de modelo. Bajo el diseño muestral probabilístico de una encuesta como la ECV, el estimador directo de una proporción se construye mediante los factores de expansión del diseño, según la Ecuación~\ref{eq:estimador-directo}, y es aproximadamente insesgado por diseño \cite{SarndalSwenssonWretman_1992}.

\begin{equation}
    \hat{y}_d = \frac{\sum_{h \in s_d} \pi_h^{-1}\, \mathbb{1}(c_h \geq k)}{\sum_{h \in s_d} \pi_h^{-1}}
    \label{eq:estimador-directo}
\end{equation}

\noindent donde:

\begin{itemize}
    \item $\hat{y}_d$: estimación directa del IPM para el departamento $d$.
    \item $s_d$: conjunto de unidades de la muestra pertenecientes al dominio $d$.
    \item $\pi_h$: probabilidad de inclusión de la unidad $h$ en la muestra.
    \item $\mathbb{1}(c_h \geq k)$: identificación del hogar como pobre multidimensional, según la Ecuación~\ref{eq:incidencia-ipm}.
\end{itemize}

La propiedad de insesgamiento por diseño es la que justifica tratar la estimación directa como punto de referencia empírico del ejercicio. Su limitación es la precisión: en dominios con muestra reducida, la varianza de muestreo crece hasta comprometer la utilidad del estimador, que es precisamente el problema que los modelos SAE buscan resolver \cite{RaoMolina_2015}.

La varianza directa se calcula como en la Ecuación~\ref{eq:varianza-directa}, como el cuadrado del error estándar reportado para cada estimación directa:

\begin{equation}
    v_d = [SE(\hat{y}_d)]^2
    \label{eq:varianza-directa}
\end{equation}

\noindent donde:

\begin{itemize}
    \item $\hat{y}_d$: IPM directo estimado para el departamento $d$.
    \item $SE(\hat{y}_d)$: error estándar de la estimación directa.
    \item $v_d$: varianza directa asociada al departamento $d$.
\end{itemize}

En la arquitectura de los modelos de área, esta varianza se trata como una cantidad fija y conocida, derivada del diseño muestral de la encuesta y no de la especificación del modelo sintético \cite{RaoMolina_2015}. Este supuesto es el que permite que la estimación directa y la información auxiliar se combinen de manera óptima, tal como se describe en la Subsección~\ref{subsec:fay-herriot}.

%--------------------------------------------------------------------------------------
\subsubsection{Matriz de covariables auxiliares y agregación de dominios}
\label{subsec:covariables}

La información auxiliar se organiza en una matriz $\mathbf{X}$ cuyas filas corresponden a las unidades territoriales y cuyas columnas corresponden a los predictores considerados. Dado que la estimación directa del IPM está disponible únicamente a escala departamental, mientras que las covariables se construyen a escala municipal, la calibración de los modelos de área requiere agregar esa información al nivel departamental.

El promedio aritmético simple resulta inadecuado para este propósito, ya que asigna el mismo peso relativo a municipios con magnitudes poblacionales muy distintas, lo que introduce un sesgo de agregación. La agregación se fundamenta entonces en promedios ponderados por la subpoblación de referencia de cada indicador, según la Ecuación~\ref{eq:agregacion-ponderada}.

\begin{equation}
    \bar{X}_{d}^{(j)} = \frac{\sum_{i \in d} p_i^{(j)}\, X_{i}^{(j)}}{\sum_{i \in d} p_i^{(j)}}
    \label{eq:agregacion-ponderada}
\end{equation}

\noindent donde:

\begin{itemize}
    \item $\bar{X}_{d}^{(j)}$: valor agregado de la covariable $j$ para el departamento $d$.
    \item $X_{i}^{(j)}$: valor de la covariable $j$ en el municipio $i$.
    \item $p_i^{(j)}$: denominador poblacional pertinente al indicador $j$ en el municipio $i$.
\end{itemize}

La elección del denominador $p_i^{(j)}$ responde a la unidad de exposición de cada fenómeno: los indicadores educativos se ponderan por la población entre 5 y 16 años, que corresponde al grupo expuesto al riesgo de rezago o inasistencia; los indicadores de mortalidad infantil se ponderan por el total de nacidos vivos, que constituye la exposición epidemiológica real; y los indicadores de infraestructura y contexto económico se ponderan por la población total municipal, dada su naturaleza de cobertura universal.

Corresponde señalar que la ponderación corrige el peso relativo de los municipios, pero no elimina la pérdida de información inherente a cualquier agregación. La varianza total de una covariable admite una descomposición entre la variación que ocurre entre departamentos y la que ocurre entre los municipios de un mismo departamento, como expresa la Ecuación~\ref{eq:descomposicion-varianza}.

\begin{equation}
    \underbrace{\sum_{d}\sum_{i \in d} (X_i - \bar{X})^2}_{\text{variación total}} =
    \underbrace{\sum_{d} n_d (\bar{X}_d - \bar{X})^2}_{\text{entre departamentos}} +
    \underbrace{\sum_{d}\sum_{i \in d} (X_i - \bar{X}_d)^2}_{\text{dentro de cada departamento}}
    \label{eq:descomposicion-varianza}
\end{equation}

Cuando la mayor parte de la variación es intradepartamental, condensar los datos en una sola medida departamental oculta los contrastes internos del territorio. Esta restricción afecta de manera desigual a las formulaciones evaluadas en este trabajo: el modelo de área utiliza únicamente los promedios departamentales $\bar{X}_d$, mientras que las formulaciones de desagregación municipal evalúan el modelo sobre la covariable específica de cada municipio $X_i$ antes de agregar sus resultados, lo que les permite conservar la dispersión interna de cada dominio.

%--------------------------------------------------------------------------------------
\subsubsection{Datos de entrenamiento y prueba}
\label{subsec:validacion}

La evaluación del desempeño en estimación para áreas pequeñas exige un esquema de validación adaptado a la estructura jerárquica de los datos. A diferencia de los problemas de aprendizaje supervisado convencionales, donde la unidad de análisis y la unidad de inferencia coinciden, en este tipo de ejercicio el parámetro de interés se calibra a nivel de dominios agregados para proyectar estimaciones sobre subáreas que no cuentan con estimación directa observable.

Bajo esta estructura, los esquemas habituales de validación cruzada por partición aleatoria de observaciones resultan inadecuados. Dividir municipios de forma aleatoria filtraría información del mismo departamento hacia los conjuntos de ajuste y de evaluación, midiendo la capacidad del modelo para reproducir patrones ya observados en lugar de su desempeño sobre dominios no observados.

La alternativa adoptada en este trabajo es la validación por exclusión de dominio completo (Leave-One-Domain-Out, LODO). En cada iteración $d = 1, \dots, D$ se excluye la totalidad de la información del dominio $d$, se reestiman los parámetros del modelo utilizando únicamente los $D-1$ dominios restantes, y se estima el indicador del dominio omitido empleando exclusivamente la componente sintética, según la Ecuación~\ref{eq:lodo}.

\begin{equation}
    \hat{\theta}_d^{(-d)} = f\!\left(\mathbf{X}_d,\; \hat{\boldsymbol{\beta}}^{(-d)}\right)
    \label{eq:lodo}
\end{equation}

\noindent donde $\hat{\boldsymbol{\beta}}^{(-d)}$ denota el vector de parámetros estimado sin la información del dominio $d$ y $f(\cdot)$ la forma funcional de la especificación evaluada. Este procedimiento reproduce el escenario operativo del estudio: estimar el indicador en unidades territoriales que no disponen de información directa.

La validez del esquema depende de que ninguna información del dominio excluido intervenga en su propia estimación. Por ello se aplican tres reglas de manera uniforme a las tres formulaciones evaluadas: el dominio excluido no recibe el ajuste de coherencia con los agregados oficiales descrito en la Subsección~\ref{subsec:no-lineal}, dado que ese procedimiento se apoya en la estimación directa que se busca reproducir; los parámetros de estandarización de las covariables se recalculan en cada iteración utilizando solo los dominios de entrenamiento; y la optimización se ejecuta con múltiples arranques para reducir el riesgo de converger a óptimos locales distintos entre iteraciones.

La estandarización de predictores mencionada consiste en una transformación $z$, definida en la Ecuación~\ref{eq:zscore}, que mitiga los problemas de mal condicionamiento numérico derivados de las diferencias de escala entre covariables.

\begin{equation}
    z_i^{(j)} = \frac{X_i^{(j)} - \mu_j^{(-d)}}{\sigma_j^{(-d)}}
    \label{eq:zscore}
\end{equation}

\noindent donde $\mu_j^{(-d)}$ y $\sigma_j^{(-d)}$ corresponden a la media y la desviación estándar de la covariable $j$ calculadas exclusivamente sobre los municipios de los dominios de entrenamiento.

%--------------------------------------------------------------------------------------
\subsubsection{Datos faltantes y técnicas de imputación}

La presencia de registros incompletos en las fuentes auxiliares admite dos tratamientos generales. El primero consiste en imputar los valores ausentes mediante algoritmos que estiman el dato faltante a partir de la información disponible en el resto de la matriz. El segundo consiste en restringir el análisis a las unidades con información completa en todas las variables requeridas, procedimiento conocido como análisis de casos completos.

La imputación resulta pertinente cuando la proporción de datos ausentes es elevada y su exclusión comprometería la representatividad del dominio, pero introduce valores que no provienen de la observación directa y cuya incertidumbre no siempre se propaga de manera explícita hacia las estimaciones finales. El análisis de casos completos evita esa fuente adicional de incertidumbre, a costa de reducir el tamaño de la muestra analítica.

En este trabajo se adopta el análisis de casos completos, aplicando un criterio de completitud común a las tres formulaciones evaluadas. La razón es de comparabilidad: al exigir que todas las formulaciones se ajusten sobre exactamente el mismo conjunto de unidades territoriales, las diferencias observadas en las métricas de desempeño resultan atribuibles a la especificación de cada modelo y no a variaciones en la muestra subyacente. Esta decisión es viable en la medida en que la proporción de registros excluidos sea marginal, condición que se verifica en la Sección~\ref{sec:resultados}.

%--------------------------------------------------------------------------------------
\subsubsection{Modelos lineales ordinarios (OLS)}
\label{subsec:ols}

La estimación por Mínimos Cuadrados Ordinarios (OLS, por sus siglas en inglés) no constituye una formulación competidora para la estimación final del IPM dentro de este diseño metodológico, sino una etapa previa de diagnóstico y selección de variables. El estimador OLS asume un error esférico y omite tanto la incertidumbre derivada del diseño muestral como la heterogeneidad no observada entre dominios, por lo que resulta insuficiente para la inferencia en áreas pequeñas \cite{RaoMolina_2015}. La forma general de los modelos OLS está dada por la Ecuación~\ref{eq:modelos-ols}:

\begin{equation}
    IPM_d = \beta_0 + \beta_1 X_{1d} + \beta_2 X_{2d} + \dots + \beta_k X_{kd} + \varepsilon_d,
    \label{eq:modelos-ols}
\end{equation}

\noindent donde:

\begin{itemize}
    \item $IPM_d$: IPM directo del departamento $d$.
    \item $\beta_0$: intercepto del modelo.
    \item $\beta_1, \beta_2, \dots, \beta_k$: coeficientes asociados a cada covariable.
    \item $X_{1d}, X_{2d}, \dots, X_{kd}$: covariables auxiliares agregadas para el departamento $d$.
    \item $\varepsilon_d$: término de error del modelo.
\end{itemize}

Pese a sus limitaciones como estimador final, esta especificación opera como filtro técnico por tres razones. La primera es la verificación de signos, que permite auditar si los coeficientes asociados a cada predictor mantienen la dirección esperada según el marco conceptual del IPM. La segunda es el diagnóstico de multicolinealidad, orientado a prevenir inestabilidad numérica en la inversión de matrices durante la estimación de los modelos jerárquicos. La redundancia entre covariables se cuantifica mediante el Factor de Inflación de Varianza, definido en la Ecuación~\ref{eq:vif}.

\begin{equation}
    VIF_j = \frac{1}{1 - R^2_j}
    \label{eq:vif}
\end{equation}

\noindent donde $R^2_j$ representa el coeficiente de determinación que resulta de regresar la covariable auxiliar $X_j$ sobre el resto de covariables del sistema. Siguiendo las convenciones econométricas, se aplica el umbral $VIF_j < 5$ para descartar la presencia de multicolinealidad severa en la matriz de diseño \cite{Greene_2018}.

La tercera razón es la detección de observaciones influyentes, que posibilita examinar la estructura de residuos estandarizados y las distancias de Cook a nivel departamental, identificando dominios atípicos antes de calibrar la varianza entre áreas del modelo de Fay-Herriot.

%--------------------------------------------------------------------------------------
\subsubsection{Modelo Fay-Herriot}
\label{subsec:fay-herriot}

El modelo de áreas de Fay y Herriot \cite{FayHerriot_1979} constituye el estándar metodológico para combinar información muestral directa con información sintética derivada de fuentes administrativas o censales a nivel de dominio agregado. A diferencia de los modelos lineales convencionales, su arquitectura se estructura en dos niveles estocásticos independientes.

El primer nivel corresponde al modelo muestral. La estimación directa observada $\hat{y}_d$ del IPM para el departamento $d$ se concibe como una medición sujeta a error de muestreo alrededor del valor poblacional verdadero no observado $\theta_d$, según la Ecuación~\ref{eq:fh-nivel1}.

\begin{equation}
    \hat{y}_d = \theta_d + e_d, \qquad e_d \sim N(0, v_d)
    \label{eq:fh-nivel1}
\end{equation}

\noindent donde $e_d$ representa la perturbación muestral y $v_d$ la varianza de muestreo definida en la Ecuación~\ref{eq:varianza-directa}, tratada como fija y conocida a partir del diseño de la encuesta.

El segundo nivel corresponde al modelo de enlace. La incidencia verdadera subyacente $\theta_d$ se relaciona de manera lineal con el vector de covariables auxiliares agregadas mediante un vector de parámetros fijos y un efecto aleatorio específico de dominio, como expresa la Ecuación~\ref{eq:fh-nivel2}.

\begin{equation}
    \theta_d = \mathbf{x}'_d \boldsymbol{\beta} + u_d, \qquad u_d \sim N(0, \sigma^2_u)
    \label{eq:fh-nivel2}
\end{equation}

\noindent donde $u_d$ captura la heterogeneidad no explicada entre departamentos, se asume independiente e idénticamente distribuido e independiente del error muestral $e_d$.

La integración de ambos niveles produce la especificación lineal mixta de la Ecuación~\ref{eq:fayherrior}, que es la forma en que habitualmente se presenta el modelo.

\begin{equation}
    \hat{y}_d = \mathbf{x}'_d \boldsymbol{\beta} + u_d + e_d
    \label{eq:fayherrior}
\end{equation}

Bajo esta estructura, el Predictor Lineal Insesgado Empírico (EBLUP, por sus siglas en inglés) para el dominio $d$ adopta la forma de una media ponderada entre la estimación directa y la estimación sintética de regresión, según la Ecuación~\ref{eq:eblup}.

\begin{equation}
    \hat{\theta}_d^{EBLUP} = \gamma_d\, \hat{y}_d + (1 - \gamma_d)\, \mathbf{x}'_d \hat{\boldsymbol{\beta}}
    \label{eq:eblup}
\end{equation}

La ponderación está gobernada por el factor de contracción definido en la Ecuación~\ref{eq:shrinkage}.

\begin{equation}
    \gamma_d = \frac{\hat{\sigma}^2_u}{\hat{\sigma}^2_u + v_d}
    \label{eq:shrinkage}
\end{equation}

El comportamiento de $\gamma_d$ refleja la adaptación del modelo a la incertidumbre del diseño. Cuando el dominio presenta un tamaño de muestra amplio y por tanto un error de muestreo bajo, $\gamma_d$ tiende a uno y la estimación se apoya casi exclusivamente en la información directa. Cuando la estimación directa es imprecisa por un error muestral elevado, $\gamma_d$ tiende a cero y la estimación se contrae hacia la componente sintética, lo que reduce sustancialmente la varianza a costa de introducir un sesgo controlado.

La varianza entre áreas $\sigma^2_u$ se estima mediante el método de Máxima Verosimilitud Restringida (REML, por sus siglas en inglés), que corrige por los grados de libertad consumidos en la estimación de los efectos fijos y evita el sesgo hacia abajo característico de la máxima verosimilitud ordinaria en muestras reducidas \cite{RaoMolina_2015}.

Una vez estimados los parámetros a escala departamental, la estimación municipal se obtiene proyectando los coeficientes sobre las covariables del municipio y sumando el efecto aleatorio del departamento correspondiente, como indica la Ecuación~\ref{eq:fh-municipal}.

\begin{equation}
    \hat{\theta}_{i,d} = \mathbf{x}'_{i,d} \hat{\boldsymbol{\beta}} + \hat{u}_d
    \label{eq:fh-municipal}
\end{equation}

Esta proyección presenta una limitación estructural relevante para el presente trabajo. El predictor lineal de la Ecuación~\ref{eq:fh-municipal} no está acotado, de modo que nada impide que la estimación resultante caiga fuera del intervalo admisible del indicador cuando las covariables municipales toman valores alejados del promedio departamental sobre el que se calibró el modelo. La corrección habitual consiste en truncar la estimación al rango válido, según la Ecuación~\ref{eq:acotamiento}, procedimiento que resuelve la inconsistencia de soporte pero no la causa que la origina.

\begin{equation}
    \hat{\theta}_{i,d}^{\,ac} = \min\!\left[100,\; \max(0,\; \hat{\theta}_{i,d})\right]
    \label{eq:acotamiento}
\end{equation}

Esta limitación no invalida al modelo, que sigue siendo una formulación metodológicamente robusta y de uso estándar en estimación para áreas pequeñas. Lo que revela es que la combinación de un predictor lineal con un indicador acotado, calibrado sobre promedios departamentales, resulta insuficiente cuando el objetivo es proyectar a una escala territorial más fina y heterogénea. Esta consideración es la que motiva las dos formulaciones alternativas que se describen a continuación.

%--------------------------------------------------------------------------------------
\subsubsection{Modelo de desagregación municipal no lineal acotada}
\label{subsec:no-lineal}

Esta formulación modifica dos elementos respecto al modelo de área: la forma funcional del predictor y el orden en que se realizan las operaciones de evaluación y agregación. Ambas modificaciones responden a una limitación documentada en la literatura: cuando el parámetro de interés está restringido a un intervalo acotado, el modelo lineal mixto resulta inapropiado, en particular para dominios cuyos valores se aproximan a los extremos del soporte, lo que ha motivado el desarrollo de modelos de muestreo no lineales con función de enlace logística para la estimación de proporciones y tasas de pobreza en áreas pequeñas \cite{Janicki_2020}.

En cuanto a la forma funcional, el predictor lineal se transforma mediante una función logística, definida en la Ecuación~\ref{eq:logistica}, que toma valores estrictamente contenidos en el intervalo abierto entre cero y uno para cualquier valor finito de su argumento.

\begin{equation}
    \pi_{i,d} = \frac{1}{1 + \exp\!\left(-\left[\mathbf{z}'_{i,d}\boldsymbol{\beta} + \alpha_d\right]\right)}
    \label{eq:logistica}
\end{equation}

\noindent donde $\mathbf{z}_{i,d}$ corresponde al vector de covariables estandarizadas del municipio $i$ del departamento $d$, según la Ecuación~\ref{eq:zscore}, y $\alpha_d$ a un desplazamiento específico del dominio. La estimación municipal respeta el rango admisible por construcción, sin requerir el truncamiento de la Ecuación~\ref{eq:acotamiento}. Esta propiedad es estructural y no depende de los valores que tomen los coeficientes ni las covariables.

En cuanto al orden de las operaciones, el modelo evalúa la función en cada municipio y pondera únicamente los resultados, en lugar de agregar primero las covariables y evaluar la función una sola vez sobre el promedio departamental. La estimación departamental se obtiene entonces según la Ecuación~\ref{eq:agregacion-nolineal}.

\begin{equation}
    \hat{\theta}_d = \sum_{i \in d} w_{i,d}\, \pi_{i,d},
    \qquad \text{con} \quad \sum_{i \in d} w_{i,d} = 1
    \label{eq:agregacion-nolineal}
\end{equation}

\noindent donde $w_{i,d}$ corresponde a la participación de la población del municipio $i$ en la población total de su departamento. La condición de suma unitaria garantiza que el agregado departamental conserve la escala e interpretación del indicador.

La distinción entre ambos órdenes de operación no es trivial. Dado que la función logística no es lineal, la desigualdad de Jensen establece que evaluar la función sobre el promedio de las covariables no equivale a promediar la función evaluada en cada municipio, relación que expresa la Ecuación~\ref{eq:jensen}.

\begin{equation}
    f\!\left(\sum_{i \in d} w_{i,d}\, \mathbf{z}_{i,d}\right) \neq \sum_{i \in d} w_{i,d}\, f(\mathbf{z}_{i,d})
    \label{eq:jensen}
\end{equation}

La consecuencia práctica es que esta formulación conserva la variabilidad interna de los municipios dentro de cada departamento, información que el modelo de área pierde al trabajar únicamente con promedios departamentales, según lo descrito en la Ecuación~\ref{eq:descomposicion-varianza}.

Los parámetros se estiman por máxima verosimilitud, incorporando tanto la varianza de muestreo del dominio como una varianza adicional propia de la formulación. Dado que la superficie de verosimilitud no es lineal y un único punto de partida puede converger a un óptimo local, la optimización se ejecuta con múltiples arranques, conservando la solución de menor valor entre todos los evaluados.

Tras el ajuste se aplica un procedimiento de coherencia con los agregados oficiales, conocido como benchmarking. Para cada dominio se determina el desplazamiento $\alpha_d$ de la Ecuación~\ref{eq:logistica} que iguala el agregado ponderado de las estimaciones municipales a la estimación directa departamental, condición expresada en la Ecuación~\ref{eq:benchmarking}.

\begin{equation}
    \sum_{i \in d} w_{i,d}\, \pi_{i,d}(\alpha_d) = \hat{y}_d
    \label{eq:benchmarking}
\end{equation}

Al tratarse de una función continua y estrictamente creciente en $\alpha_d$, esta ecuación admite solución única para cada dominio. La literatura identifica este tipo de ajuste como deseable por dos motivos: garantiza la coherencia institucional entre las estimaciones en áreas pequeñas y las cifras oficiales publicadas en el nivel de agregación superior, cuya mayor muestra las hace más confiables, y ofrece además cierta protección frente a errores de especificación del modelo \cite{DattaGhoshSteortsMaples_2011}.

%--------------------------------------------------------------------------------------
\subsubsection{Modelo de desagregación espacial CAR}
\label{subsec:car}

Esta formulación extiende el modelo de desagregación no lineal incorporando un componente espacial a escala municipal bajo una estructura Condicional Autorregresiva (CAR, por sus siglas en inglés), que modela explícitamente la dependencia entre unidades vecinas. La especificación condicional autorregresiva fue formalizada por Besag \cite{Besag_1974}, quien mostró que la distribución conjunta de un conjunto de variables espacialmente interactuantes puede construirse a partir de sus distribuciones condicionales respecto a la vecindad. La incorporación de efectos de área espacialmente correlacionados en modelos de estimación para áreas pequeñas ha documentado reducciones tanto en la varianza como en el sesgo del predictor frente a su versión sin componente espacial \cite{PratesiSalvati_2008, PratesiSalvati_2019}.

La estructura CAR especifica la distribución conjunta de los efectos espaciales a partir de la matriz de adyacencia del territorio, según la Ecuación~\ref{eq:car}.

\begin{equation}
    \boldsymbol{\phi} \sim N\!\left(\mathbf{0},\; \tau^{-1}(\mathbf{D} - \rho \mathbf{W})^{-1}\right)
    \label{eq:car}
\end{equation}

\noindent donde:

\begin{itemize}
    \item $\boldsymbol{\phi}$: vector de efectos espaciales municipales.
    \item $\mathbf{W}$: matriz de adyacencia, con $w_{ij}=1$ si los municipios $i$ y $j$ son vecinos y cero en caso contrario.
    \item $\mathbf{D}$: matriz diagonal cuyo elemento $d_{ii}$ corresponde al número de vecinos del municipio $i$.
    \item $\rho$: parámetro que gobierna la intensidad de la dependencia espacial.
    \item $\tau$: parámetro de precisión del componente espacial.
\end{itemize}

El predictor municipal combina entonces la componente fija de covariables con el efecto espacial, antes de aplicar la transformación logística de la Ecuación~\ref{eq:logistica}, como indica la Ecuación~\ref{eq:car-predictor}.

\begin{equation}
    \pi_{i,d} = \frac{1}{1 + \exp\!\left(-\left[\mathbf{z}'_{i,d}\boldsymbol{\beta} + \phi_i + \alpha_d\right]\right)}
    \label{eq:car-predictor}
\end{equation}

La estimación de un efecto aleatorio independiente por municipio resulta inviable cuando solo se dispone de un número reducido de dominios observados para la calibración. Por ello el componente espacial se representa mediante una versión de rango reducido, que proyecta $\boldsymbol{\phi}$ sobre los primeros $K$ autovectores de la matriz de vecindad normalizada, según la Ecuación~\ref{eq:rango-reducido}.

\begin{equation}
    \boldsymbol{\phi} \approx \mathbf{U}_K \boldsymbol{\delta}, \qquad \boldsymbol{\delta} \in \mathbb{R}^{K}, \quad K \ll n
    \label{eq:rango-reducido}
\end{equation}

\noindent donde $\mathbf{U}_K$ contiene los $K$ autovectores asociados a los autovalores de mayor magnitud, que corresponden a los patrones espaciales de mayor escala. Esta reducción conserva la estructura territorial dominante al tiempo que limita el número de parámetros a estimar.

Por la misma limitación de identificabilidad, el parámetro $\rho$ no se estima libremente sino que se fija a priori, evaluando su efecto sobre una grilla de valores como análisis de sensibilidad. La agregación departamental y el procedimiento de benchmarking operan igual que en la formulación no lineal, según las Ecuaciones~\ref{eq:agregacion-nolineal} y~\ref{eq:benchmarking}.

La construcción de la matriz $\mathbf{W}$ requiere una definición operativa de vecindad. El criterio de contigüidad tipo Reina considera vecinos a dos municipios que comparten al menos un punto o segmento fronterizo. La discontinuidad geográfica del territorio colombiano, con dominios insulares y municipios sin contigüidad terrestre, genera filas nulas en la matriz de adyacencia que producen matrices singulares en la Ecuación~\ref{eq:car}. Para resolverlo se adopta una regla de conexidad forzada basada en la distancia a los vecinos más cercanos, que garantiza que todo municipio cuente con al menos un vínculo en el grafo espacial.

Esta condición tiene una implicación interpretativa que conviene explicitar. En las formulaciones sin componente espacial, la estimación de un dominio insular se deriva íntegramente de sus covariables locales. En la formulación CAR, al no existir contigüidad terrestre, el componente autorregresivo no transfiere suavizamiento espacial desde el resto del territorio, de modo que la estimación sobre esas unidades se fundamenta esencialmente en la componente sintética ajustada.

%--------------------------------------------------------------------------------------
\subsubsection{Métricas de desempeño}
\label{subsec:metricas}

La evaluación comparativa de las formulaciones se fundamenta en un esquema de métricas complementarias, siguiendo las recomendaciones habituales para la validación de modelos de estimación en áreas pequeñas \cite{RaoMolina_2015, moura2021}. Ningún indicador individual captura la totalidad de las dimensiones del error en dominios no observados, por lo que la literatura sugiere evaluar simultáneamente la magnitud absoluta de la desviación, la sensibilidad a valores atípicos y la preservación de la estructura ordinal de los dominios. Las tres métricas que se describen a continuación se calculan sobre las estimaciones producidas bajo el esquema de validación por exclusión de dominio de la Subsección~\ref{subsec:validacion}.

La Raíz del Error Cuadrático Medio (RMSE, por sus siglas en inglés) evalúa la precisión global en la escala original del indicador, según la Ecuación~\ref{eq:rmse}.

\begin{equation}
    RMSE = \sqrt{\frac{1}{D} \sum_{d=1}^{D} \left(\hat{y}_d - \hat{\theta}_d\right)^2}
    \label{eq:rmse}
\end{equation}

Debido a la elevación al cuadrado del residuo, el RMSE penaliza con mayor severidad los errores de gran magnitud. En el contexto de la política social, un RMSE elevado advierte sobre la presencia de dominios donde el modelo sintético falla de manera acentuada, aspecto crítico al evaluar el riesgo de subestimación de la pobreza en áreas extremas.

El Error Absoluto Medio (MAE, por sus siglas en inglés) mide la magnitud promedio del error sin ponderación cuadrática, según la Ecuación~\ref{eq:mae}.

\begin{equation}
    MAE = \frac{1}{D} \sum_{d=1}^{D} \left| \hat{y}_d - \hat{\theta}_d \right|
    \label{eq:mae}
\end{equation}

Al ponderar todos los residuos de forma lineal, el MAE ofrece una medida de la exactitud típica del modelo, menos sensible a valores atípicos o a variaciones en las colas de la distribución.

El coeficiente de correlación de Pearson evalúa el grado de asociación lineal y la preservación del ordenamiento relativo entre la estimación directa y la del modelo, según la Ecuación~\ref{eq:correlacion}.

\begin{equation}
    r = \frac{\sum_{d=1}^{D} \left(\hat{y}_d - \bar{y}\right)\left(\hat{\theta}_d - \bar{\theta}\right)}
    {\sqrt{\sum_{d=1}^{D} \left(\hat{y}_d - \bar{y}\right)^2} \; \sqrt{\sum_{d=1}^{D} \left(\hat{\theta}_d - \bar{\theta}\right)^2}}
    \label{eq:correlacion}
\end{equation}

En ejercicios de focalización territorial esta métrica adquiere relevancia estratégica: un modelo con un MAE aceptable pero con baja correlación alteraría la priorización de los dominios, lo que conduciría a decisiones subóptimas en la asignación de recursos públicos.

La adopción conjunta de las tres métricas permite un diagnóstico diferencial que ninguna de ellas ofrece por separado. La discrepancia entre un MAE bajo y un RMSE alto identifica escenarios de buen ajuste generalizado con fallas puntuales concentradas en dominios específicos. La combinación de alta correlación con métricas de error controladas confirma que la formulación no solo aproxima la magnitud del indicador, sino que además conserva la jerarquía territorial de las necesidades.

Para la comparación entre especificaciones ajustadas sobre el mismo conjunto de dominios se emplean adicionalmente los criterios de información de Akaike y Bayesiano, definidos en las Ecuaciones~\ref{eq:aic} y~\ref{eq:bic}, que penalizan la verosimilitud alcanzada en función del número de parámetros estimados y favorecen especificaciones con mejor balance entre ajuste y parsimonia.

\begin{equation}
    AIC = -2\log L + 2p
    \label{eq:aic}
\end{equation}

\begin{equation}
    BIC = -2\log L + p \log D
    \label{eq:bic}
\end{equation}

\noindent donde $L$ es la verosimilitud del modelo ajustado, $p$ el número de parámetros y $D$ el número de dominios.

Finalmente, la medida de precisión del estimador de área corresponde a su error cuadrático medio. A diferencia de la varianza de muestreo del estimador directo, que se conoce por diseño, el error cuadrático medio del EBLUP no admite una expresión cerrada exacta, ya que debe incorporar la incertidumbre adicional que introduce la estimación de la varianza entre áreas. La aproximación estándar descompone esta cantidad en tres términos: la contribución del predictor con parámetros conocidos, la derivada de la estimación de los coeficientes fijos y la asociada a la estimación de la varianza del efecto aleatorio \cite{PrasadRao_1990}. A partir de esa aproximación, el error estándar reportado se obtiene según la Ecuación~\ref{eq:se}.

\begin{equation}
    SE_{FH,d} = \sqrt{MSE_d}
    \label{eq:se}
\end{equation}

%--------------------------------------------------------------------------------------
\subsubsection{Índice de Moran}
\label{subsec:moran}

El estadístico global de Moran \cite{Moran_1950, Anselin_1988} cuantifica el grado de autocorrelación espacial de una variable distribuida sobre un conjunto de unidades territoriales. En este trabajo se aplica sobre los residuos de las especificaciones exploratorias sin componente espacial, durante la etapa de análisis exploratorio descrita en la Sección~\ref{sec:metodologia}, con el fin de determinar si existe estructura territorial no capturada por las covariables. Su formulación corresponde a la Ecuación~\ref{eq:i-moran}.

\begin{equation}
    I = \frac{D}{S_0} \cdot \frac{\sum_i \sum_j w_{ij}(e_i - \bar{e})(e_j - \bar{e})}{\sum_i (e_i - \bar{e})^2}
    \label{eq:i-moran}
\end{equation}

\noindent donde:

\begin{itemize}
    \item $I$: estadístico global de Moran.
    \item $D$: número de dominios evaluados.
    \item $w_{ij}$: peso espacial entre los dominios $i$ y $j$.
    \item $e_i$ y $e_j$: residuos del modelo en los dominios $i$ y $j$.
    \item $\bar{e}$: promedio de los residuos.
    \item $S_0$: suma total de los pesos espaciales, $S_0 = \sum_i \sum_j w_{ij}$.
\end{itemize}

El estadístico se encuentra acotado aproximadamente en el intervalo $[-1, 1]$. Un valor significativamente positivo indica que dominios geográficamente próximos presentan perturbaciones similares, lo que revela una estructura de dependencia territorial no capturada por el conjunto de covariables fijas.

Para evitar que la conclusión dependa de una definición discrecional del entorno geográfico, la inferencia se realiza mediante pruebas de permutación condicional, evaluando la consistencia del estadístico bajo distintos esquemas de vecindad. La detección de autocorrelación espacial significativa en esta etapa constituye la justificación empírica para formular el modelo espacial descrito en la Subsección~\ref{subsec:car}, que incorpora dicha estructura de manera explícita en la matriz de covarianza de los efectos aleatorios.

%--------------------------------------------------------------------------------------
\subsubsection{Supuestos distribucionales y diagnósticos de validez}
\label{subsec:diagnosticos}

La validez de las estimaciones depende del cumplimiento razonable de los supuestos distribucionales de cada formulación. Los diagnósticos que se enumeran a continuación se aplican en la Sección~\ref{sec:metodologia} y sus resultados se reportan en la Sección~\ref{sec:resultados}.

La normalidad de los residuos estandarizados se evalúa mediante la prueba de Shapiro-Wilk, cuya hipótesis nula establece que la muestra proviene de una distribución normal. La homocedasticidad se examina mediante la prueba de Breusch-Pagan, que contrasta la hipótesis nula de varianza constante del error frente a la alternativa de que esta dependa de las covariables. La adecuación de la forma funcional lineal se verifica mediante la prueba RESET, que evalúa si potencias del predictor ajustado aportan capacidad explicativa adicional, lo que señalaría una especificación funcional incompleta \cite{Greene_2018}.

En las formulaciones estimadas por optimización numérica se verifica adicionalmente la identificación local del modelo, comprobando que la matriz Hessiana evaluada en la solución sea definida positiva, y la estabilidad de la solución frente a múltiples arranques del algoritmo.

Estos diagnósticos se interpretan como elementos que documentan el comportamiento del ajuste y delimitan el alcance de las conclusiones, y no como reglas mecánicas de descarte. Una advertencia en alguno de ellos identifica una condición que debe considerarse al interpretar los resultados, sin que ello invalide por sí misma una formulación cuyo sustento metodológico es reconocido en la literatura del área.